\documentclass[11pt]{article}
\usepackage[margin=1in]{geometry}
\usepackage{amsmath,amssymb,amsthm}
\usepackage{booktabs,longtable,array}
\usepackage[T1]{fontenc}
\usepackage{enumitem}

\newtheorem{definition}{Definition}
\newtheorem{proposition}{Proposition}
\newtheorem{theorem}{Theorem}
\newtheorem{corollary}{Corollary}
\newtheorem{remark}{Remark}
\newcommand{\Ftwo}{\mathbb F_2}
\newcommand{\Bottom}{\mathsf{Bottom}_{\mathrm{src}}}

\title{V21 P1-T02 Proposal-Only EqSrc Canonical-Choice\\
Candidate-Family and Lineage Map}
\author{Ontology Formalizer control packet RT-20260720-010}
\date{2026-07-20}

\begin{document}
\maketitle

\section{Control status and exact question}

This artifact executes only v21 work item \texttt{P1-T02}. It assigns stable
family identities and immutable candidate identities to the recent
intrinsic-discriminator, cycle-boundary, orientation-torsor,
rooted-partition, and graded-orbit candidate chain. It records exact
construction, audit, stress, repair, local-freeze, and route-successor
lineage without changing any candidate's scientific status.

Every candidate remains \texttt{proposal-only}, \texttt{draft/control}, and
\texttt{source-extension data}. Current ontology does not derive the
decisive source-side selector. The adoption status is
\texttt{blocked\_adoption\_open\_continuation}. This packet does not repair,
audit, stress, freeze, adopt, reject, or promote a candidate. P1-T03 owns the
separate family-freeze or materially-distinct-route decision.

\section{Stable family and candidate identities}

\begin{longtable}{p{0.25\linewidth}p{0.22\linewidth}p{0.43\linewidth}}
\toprule
Family ID & Candidate IDs & Exact scope \\
\midrule
\endfirsthead
\toprule
Family ID & Candidate IDs & Exact scope \\
\midrule
\endhead
\texttt{EQSRC-CANONICAL-CHOICE-INTRINSIC-DISCRIMINATOR} &
\texttt{EQSRC-IDISC-V1}, \texttt{-V2}, \texttt{-V3} &
Finite two-truncated $\Ftwo$ chain packages, homology discriminator,
boundary subspace, kernel-pair relation, and boundary translations. \\
\texttt{EQSRC-CANONICAL-CHOICE-CYCLE-BOUNDARY-LINE} &
\texttt{EQSRC-CYCLE-BOUNDARY-V1} &
Odd oriented mark, marked automorphisms, invariant state space, odd-mark
functional, and its kernel line. \\
\texttt{EQSRC-CANONICAL-CHOICE-ORIENTATION-TORSOR} &
\texttt{EQSRC-ORIENTATION-TORSOR-V1} &
Unordered both-odd bipartition, free $C_2$ orientation torsor, equivariant
fibre lines, and associated state--relation object. \\
\texttt{EQSRC-CANONICAL-CHOICE-ROOTED-PARTITION} &
\texttt{EQSRC-ROOTED-PARTITION-V1} &
Transition certificates, event projection, singleton root, path parity,
ordered partition, boundary line, and kernel-pair relation. \\
\texttt{EQSRC-CANONICAL-CHOICE-GRADED-ORBIT-ROOT} &
\texttt{EQSRC-GRADED-ORBIT-ROOT-V1} &
Ordered-monoid action, full component, reachability covers, global minimum,
rank parity, boundary line, and kernel-pair relation. \\
\bottomrule
\end{longtable}

The machine-readable inventory stores, for every candidate, the exact source
artifact SHA-256, a canonical hash of its normalized theorem-statement
signature, a canonical hash of its assumption signature, and the current
ontology anchor hash. These hashes are identity and provenance evidence; they
are not proof or adoption authority.

\begin{definition}[Exact candidate identity and duplication]
For an inventory entry $c$, let
\[
 I(c)=\bigl(\operatorname{id}(c),F(c),v(c),H_{\rm art}(c),
 H_{\rm thm}(c),H_{\rm asm}(c),H_{\rm ont}(c)\bigr).
\]
Two inventory entries are \emph{exact duplicates} only if their artifact,
theorem-statement, assumption, and ontology-anchor hashes agree. Similar
names, isomorphic finite instances, or a shared kernel-pair codomain are not
identity criteria.
\end{definition}

\section{Exact lineage graph}

The graph contains two within-family repair edges
\[
 \texttt{IDISC-V1}\longrightarrow\texttt{IDISC-V2}
 \longrightarrow\texttt{IDISC-V3}
\]
and four distinct-family route-successor edges
\[
 \texttt{IDISC-V3}\longrightarrow\texttt{CYCLE-BOUNDARY-V1}
 \longrightarrow\texttt{ORIENTATION-TORSOR-V1}
 \longrightarrow\texttt{ROOTED-PARTITION-V1}
 \longrightarrow\texttt{GRADED-ORBIT-ROOT-V1}.
\]
The first two edges are semantic theorem-scope repairs and supersede the
earlier variant inside one family. The last four edges are materially
distinct route successors after an exact-cycle local freeze. They do not
assert cross-family identity or semantic supersession.

\begin{proposition}[Acyclicity of the exact candidate graph]
The candidate lineage graph recorded by P1-T02 is acyclic.
\end{proposition}

\begin{proof}
Assign construction ordinal
\[
 27<29<31<35<39<43<47
\]
to the seven candidates in the displayed order. Every registered edge
strictly increases this ordinal. A directed cycle would give a finite strict
chain returning to its starting ordinal, which is impossible.
\end{proof}

\begin{proposition}[No exact duplicates in the bounded inventory]
No two distinct candidates, and no two distinct families, in the bounded
P1-T02 inventory are exact duplicates.
\end{proposition}

\begin{proof}
The task-local validator recomputes every artifact, theorem-statement,
assumption, and ontology-anchor hash. Intrinsic-discriminator v1, v2, and v3
share a core relation constructor but have different theorem and assumption
signatures after their substantive repairs. Every cross-family pair has a
different source domain or selector primitive and therefore a different
assumption signature. Hence no distinct pair satisfies the exact-duplicate
criterion.
\end{proof}

\section{Inventory-scoped selector criterion}

\begin{definition}[Line-selected kernel pair]
Let $A$ be an $\Ftwo$-vector space and let a selector value $s$ determine a
subspace $B_s\subseteq A$. Define
\[
 x\,R_s\,y\quad\Longleftrightarrow\quad x-y\in B_s.
\]
This definition is algebraic. It does not state where $s$ comes from or that
any selector is physically admissible.
\end{definition}

\begin{theorem}[Distinct selected lines give distinct relations]
If $B_s\ne B_t$, then $R_s\ne R_t$. Consequently, a candidate's admission
axioms do not determine a unique relation whenever they admit two selector
values with distinct selected subspaces.
\end{theorem}

\begin{proof}
Choose $b$ in the symmetric difference of $B_s$ and $B_t$; without loss of
generality $b\in B_s\setminus B_t$. Then $0\,R_s\,b$, while
$0\not\mathrel{R_t}b$. Thus the relations differ. The second statement is
the direct application of this witness to the admitted selector values.
\end{proof}

\begin{remark}[Exact scope]
This theorem is not the general source-natural selector theorem reserved for
P2. It is a small algebraic criterion used only to normalize the already
tracked finite evidence. It does not show that every future source extension
admits multiple selector values.
\end{remark}

\section{Tracked finite instantiations and repeated obstruction}

\begin{longtable}{p{0.25\linewidth}p{0.65\linewidth}}
\toprule
Family & Tracked bounded evidence \\
\midrule
\endfirsthead
\toprule
Family & Tracked bounded evidence \\
\midrule
\endhead
Intrinsic discriminator & Exact v3 stress retains four valid differentials
on fixed carriers that induce different relations; current ontology does not
select the differential. \\
Cycle boundary & The bounded stress exposes multiple line choices under the
broader symmetry and an orientation-descent and independent-selection
obstruction. \\
Orientation torsor & The bounded stress exposes multiple equivariant
descended lines and no source-derived physical line or variation selector. \\
Rooted partition & Among the enumerated four-token cases, 265 admitted
transition graphs support two root-induced relations, and all 182 distinct
ordered odd-block changes change the relation. \\
Graded orbit & Every nontrivial ordered finite partition is realizable by
freely supplied action data; the local record also fails to certify its
ambient component completion. \\
\bottomrule
\end{longtable}

The shared obstruction identity is
\texttt{OB-EQSRC-CANONICAL-CHOICE-SOURCE-PROVENANCE-RECURS-001}.
It says only that the current ontology does not derive the decisive
differential, oriented mark, partition line, root, ordered action, component,
physical morphism, or physical variation selector in these tracked families.
The finite algebra survives after the data are supplied. No global no-go,
future-source-extension impossibility, or program-level rejection follows.

\section{Representation versus genuine assumption change}

No registered edge is classified as representation-only. The v1-to-v2 and
v2-to-v3 intrinsic-discriminator edges preserve the kernel-pair core but make
substantive semantic and typed-theorem repairs. Each cross-family edge changes
the supplied source primitives: differential; oriented mark; unordered
partition and torsor; transition, projection, and root; or ordered action,
full component, and grading. The task-local assumption-delta table records
these changes exactly.

\section{Source-extension and no-target classification}

The whole candidate objects are proposal-only new-primitive candidates under
\texttt{blocked\_adoption\_open\_continuation}. Their conditional algebra is
not reclassified as derived from current ontology. Prior audits, stresses,
selector decisions, and this lineage map are status-boundary evidence only.

No target topology, target smooth atlas, target metric, proper time, detector
semantics, benchmark success, generated derivative, registry row, validator,
role, handoff, file order, or task-process status is used as a source premise.
A validator PASS establishes structural consistency only. Absence of explicit
target or process import does not establish source provenance, physical
admissibility, or derivation from current ontology.

\section{Distance-to-GR and freeze boundary}

\begin{center}
\begin{tabular}{p{0.35\linewidth}p{0.5\linewidth}}
\toprule
Object & Status after P1-T02 \\
\midrule
General $\mathsf{EqSrc}$ & Undischarged; no adopted source law. \\
Distance-to-GR ledger & Unchanged. \\
Metric-use ledger & Unchanged. \\
Prior exact candidate cycles & Existing local freezes preserved exactly. \\
Active graded-orbit family & \texttt{freeze\_review\_required}; no decision here. \\
P1-T03 & Dependency-next and not executed by this packet. \\
Canonical ontology and downstream GR & No edit, adoption, derivation, or promotion authorized. \\
\bottomrule
\end{tabular}
\end{center}

\section{Conclusion}

P1-T02 creates one auditable family inventory, one acyclic candidate lineage
graph, and one assumption-delta table. It identifies five stable families,
seven immutable candidates, no exact duplicates, two within-family repairs,
and four distinct-family route-successor edges. The recurrent scoped
obstruction is preserved without overread. P1-T03 must decide whether to
freeze the families at family level or permit one materially distinct theorem
route.

\end{document}
